\textbf{Temporal attention and occlusion.} The [CLS] attention of the last
block is spread almost evenly over the twelve segments (0.07--0.09 per
segment, against 0.083 for uniform attention), with class-specific emphasis: the
middle of the window for jumping jacks, early segments for squats, and the
first segment for walking (Fig.~\ref{fig:xai}). Removing a single segment
(occlusion) lowers the true-class probability by less than one percentage point
on average. The model therefore integrates evidence over the whole window and
does not depend on any single 170-ms segment, which is consistent with its
tolerance to frame loss. Attention and occlusion importance agree only
moderately (Pearson $r=\xaiCorr$ over the class$\times$segment cells), so we
treat attention as a coarse indication and not as a faithful explanation.
\textbf{Point saliency.} Gradient saliency concentrates on the returns of the
moving limbs and torso and is low on isolated returns far from the body,
i.e.\ on the residual clutter that survives gating. \textbf{Embeddings.} In a
t-SNE of RadHAR test embeddings the classes form compact, well-separated
clusters (silhouette \silSTPTLiteKDrobust{} in the model's embedding space
versus \silPointNet{} for PointNet). \textbf{Errors.} The remaining RadHAR
errors are mostly windows predicted as \emph{jump} (precision 91\%), which
shares the strong vertical Doppler pattern of jumping jacks and squats. On
MiliPoint, errors concentrate in groups of kinematically similar actions.
